% Original Peirce contribution to planned 176.2; distinct from Psillos 2011.
\chapter{Peirce: Explanatory Conjecture and the Conduct of Inquiry}
\label{sec:p16-peirce-abduction-and-inquiry}

An observation can prompt a conjecture without establishing it. A conjecture
can explain an observation without being the only explanation. Testing a
conjecture requires more than redescribing the observation which suggested
it. Peirce's accounts give these differences a place in the conduct of
inquiry, but they do not present one unchanging mathematical relation
called abduction.

Two published primary presentations are distinguished here. In the sixth
\emph{Illustrations of the Logic of Science} paper, published in 1878,
Peirce classifies deduction, induction and \emph{hypothesis} through the
roles of rule, case and result. In \emph{A Neglected Argument for the
Reality of God}, published in 1908, he describes \emph{retroduction},
deduction and induction as stages of inquiry. The latter's theological
application, account of instinct and express limitations are part of
the setting in which its methodology is proposed. They must not be
removed to make the text look like an already proved modern algorithm
\parencite[secs.~I--V]{Peirce1878DIHPrimary}
\parencite[secs.~I--V]{Peirce1908NeglectedPrimary}.

The comparison concerns these two presentations, not every intervening
manuscript or lecture. In particular, the familiar modern finite
explanation relation in \cref{sec:p16-finite-abduction} is separately
attributed to Eiter and Gottlob. It is not Peirce's original system under
another name.

\section{Rule, case and result in the 1878 classification}
\label{sec:p16-peirce-1878-roles}

Peirce opens by connecting the testing of arguments with their
classification. Barbara, the typical first-figure syllogism, represents
the application of a general rule to a case. His mortality example
distinguishes a rule about men, the classification of Enoch and Elijah,
and the resulting conclusion. Its function is to exhibit argument roles,
not to establish the example's historical or theological premises.

He also suggests that arguments can be put into Barbara's form if its
copula is allowed to mean representation for the purpose of the argument,
rather than literal inclusion alone. But he explicitly declines to give
the general proof. More importantly, he says that a common representational
form need not exhibit the distinctive operation of each inference. Calling
a sample a representative of a population does not prove that it represents
that population adequately \parencite[sec.~I]{Peirce1878DIHPrimary}.

His quadrangle example illustrates a related background issue. The
reconstruction of the inference to some figures which are not triangles
contains a premise that some figures are quadrangles. The existential
commitment is not obtained merely by writing a universal classification.
The example presupposes an applicable case; it is not a theorem that
universal statements always establish the existence of their subjects.

\begin{hbdefinition}{def:p16-peirce-role-distinction}{The three argument roles}
In the 1878 presentation, a \emph{rule} supplies a general classification
or connection; a \emph{case} places specified objects under it; a
\emph{result} states the character which the connection assigns to that
case. Deduction applies the rule to the case. Induction proposes a rule
from a case and its observed result. Hypothesis proposes a case from a
rule and the observed result.

These are Peirce's roles in this presentation, not three interchangeable
directions of a single truth-preserving consequence relation. The latter
two proposals may be wrong while their starting observations are right.
\end{hbdefinition}

The statistical version does not change this distinction. If the
proportion of white beans in a bag is already supplied and the drawing
procedure has the required representative character, a probability
assessment about a drawn bean applies that supplied rule. Estimating the
bag's unknown proportion from a drawn handful instead proposes the rule.
Peirce describes the relevant drawing through a long-run relative-frequency
condition. That description does not itself establish a physical sampling
mechanism, independence, or a modern convergence theorem.

\section{The beans and what each conclusion adds}
\label{sec:p16-peirce-beans}

\begin{hbexample}{def:p16-peirce-beans}{One setting, three inferential tasks}
\emph{Deduction.} Suppose every bean from a specified bag is white and
these beans came from that bag. Applying the rule gives their whiteness.
The source of the beans is already a premise.

\emph{Induction.} Suppose these beans came from that bag, their color is
observed, and the handful is taken as an appropriate sample. Proposing
that all beans in the bag are white extends the observation to beans
not examined. Peirce also gives the proportional version, in which a
two-thirds-white sample suggests approximately that proportion in the
bag. Neither observation excludes a differently colored unexamined bean.

\emph{Hypothesis.} Suppose these beans are white and a bag contains only
white beans. Proposing that the beans came from that bag identifies an
unobserved source. Peirce's room contains several bags, and the white
handful lies on the table. Finding a suitable bag supplies a possible
explanation, not a deduction of the handful's provenance.
\end{hbexample}

A concrete failure case preserves the rule and result of the third
argument: a second bag also contains white beans, and the handful came
from that second bag. The first bag's rule remains true, and the handful
remains white, while the proposed case is false. This is why the
explanatory direction is not the deductive direction reversed as a
truth-preserving operation. A more discriminating observation might
favor one origin over another, but the three-role pattern alone supplies
neither that observation nor its reliability.

Peirce describes hypothesis as sometimes only a surmise, not a settled
belief. His other introductory examples have the same source-searching
character. A lavishly escorted rider suggests the governor because of
the honors shown him; inland fish fossils suggest an earlier sea; and
documents and monuments suggest Napoleon's existence. The explanatory
connection in each case uses background understanding. Escort conventions,
formation of fossils and production of documents are not consequences
of the bare observed descriptions alone
\parencite[sec.~I]{Peirce1878DIHPrimary}.

The examples range from weak conjecture to familiar historical assurance.
That range does not license assigning all hypotheses one probability,
or treating a physically unobserved object as necessarily unknowable.
The source distinguishes how something is inferred from whether one
can directly observe the proposed fact. Peirce also compares a tentative
hypothesis with belief in yesterday's events prompted by the feeling of
remembering them, proposing a difference of degree rather than a wholly
different inferential kind. This is his philosophical comparison, not
a measured account of memory reliability.

\section{Indirect moods and the limits of inversion}
\label{sec:p16-peirce-indirect}

Section~II compares the two synthetic inferences with indirect
syllogistic moods. If a necessary consequence is denied while a rule
is retained, a proposed case can have to be denied; if the case is
retained, the unrestricted rule can have to be denied. Peirce uses
Baroco and Bocardo, with the mortality illustration, to make these
roles visible. The deductive versions retain the necessary connection.

He next replaces that connection with a probable one. A bag mostly
containing white beans would normally suggest a mostly white handful.
A contrary-looking handful can prompt either a different origin or a
different estimate of the bag's composition. This comparison draws
attention to competing explanations of the result rather than merely
to agreement with one favored explanation.

Peirce then qualifies his own comparison. Induction is genuinely the
proposal of a rule, not only the denial of another rule. Hypothesis is
genuinely a classification under a proposed origin or class, not only
a denial of another origin. The indirect description can become
artificial. He expressly distinguishes necessary contraposition from
probable inference: a probable consequence does not have the same
unqualified contrapositive as a necessary consequence
\parencite[sec.~II]{Peirce1878DIHPrimary}.

For example, a minority-colored handful remains possible when most beans
in the original bag are white. Without a specified sampling distribution,
sample size and competing origin model, calling that handful unusual
does not determine a numerical assessment of its origin. The historical
pattern is useful for locating a question; it is not a ready-made
likelihood ratio, posterior probability or significance test.

\section{Explanation is not simply more resemblance}
\label{sec:p16-peirce-resemblance}

Peirce's torn-paper example adds a more distinctive character. An
anonymous writing occupies one torn piece; a piece in a suspected
author's desk matches its irregular edge, and only that person had
access to the desk. Peirce takes the unlikely accidental fit to favor
the authorship conjecture. But he distinguishes this conclusion from
predicting further matching irregularities. The latter would extend
observed characters; the former connects a physical match with a
different fact, the proposed writer's identity
\parencite[sec.~III]{Peirce1878DIHPrimary}.

He also warns that characters are not enumerated in the same straightforward
way as sampled objects, and that they occur in categories. Several described
matches therefore need not be several independently informative observations.
The source gives no feature-independence model which would justify treating
each additional description as a separate multiplier of support.

The assumptions matter. An improbable accidental match is not by itself
a probability of authorship. Even if the two pieces belong together,
someone else could have written on the sheet before the suspected
person acquired it. That possibility can respect exclusive access to
the desk. Authorship therefore needs a provenance connection in addition
to a match. This is an editorial separating case, not a report about
an actual investigation or legal evidence standard.

The solar-spectrum example combines an explanatory conjecture with
additional knowledge about distinctive spectral lines. Peirce says that
this knowledge strengthens the metal-presence hypothesis. It is not the
same as counting any chosen similarities. His solar-myth/Napoleon parody
illustrates the danger of selecting resemblances freely after choosing
what one wants to explain. Ingeniously chosen correspondences can be
made to favor a fanciful story.

Peirce rejects making a blanket uniformity of Nature the foundation of
hypothesis, while allowing that knowledge of particular uniformities can
strengthen a particular inference. The solar lines illustrate the latter
role. His rejection is not a denial that regularities exist, and the
reference to warnings in an earlier paper does not supply that paper's
complete justification of induction here. Neither naming uniformity nor
rejecting its blanket use removes the need to defend the actual connection.

The paper announces, but omits, a general demonstration concerning
resemblance. No such omitted demonstration is supplied merely by these
examples. The bounded lesson is that unrestricted selection of matching
characters does not establish a general support relation. Resemblance
can still be informative under a warranted mechanism or observation
model. The mechanism and selection policy are the missing work, not
something made irrelevant by the word hypothesis.

\begin{hbmethod}{def:p16-peirce-1878-tests}{The 1878 safeguards for testing a hypothesis}
Peirce requires that the hypothesis be posed distinctly as a question
before the observations which test it. Predictions are then to be
examined rather than merely repeating the observations which suggested
the conjecture. The respects compared must not be selected because
the hypothesis is already known to work well on them; he describes
their selection as random. Failures as well as successes must be recorded.

These are methodological requirements about inquiry. Their implementation
needs an actual set of possible tests, a selection procedure and an honest
record. The text does not supply a universal randomization design, error
bound or guarantee that every compliant hypothesis is true.
\end{hbmethod}

His attacks on obstinacy, authority and partisan attempts to reach a
foregone conclusion defend the purpose of learning how matters stand.
Their force is not established by his assertion that proofs could be
given elsewhere. Nor does his polemic about legal procedure supply a
present account of law. Mill's four methods are named as valuable
additional contrivances; their complete arguments are not contained in
this paper \parencite[sec.~III]{Peirce1878DIHPrimary}.

The distinction between discovery and test is not a prohibition on
learning a hypothesis from data. It prevents giving the same selected
fit two different evidential jobs without accounting for the selection.
An observation can generate a conjecture and enter a later assessment,
but its prior use must remain visible when that assessment is defended.

\section{Empirical formulas, mechanisms and mixed reasoning}
\label{sec:p16-peirce-formula}

Peirce admits border cases between his classifications. In Section~IV,
observed changes under changing circumstances can suggest a relation
which looks more inventive than a simple classification. His example
fits a cubic expression to the relative volume of a fixed mass of water
at different temperatures.

\begin{hbexample}{def:p16-peirce-empirical-formula}{The cubic interpolation example and its scope}
In the source's notation, \(v\) is relative volume, \(t\) temperature,
and \(a,b,c\) are numerical coefficients to be fitted. The proposed
expression is
\[
                         v=1+at+bt^2+ct^3.
\]
The first observations suggest the form; observations at other temperatures
are then examined. The proposed extension concerns the temperature range
from which the observations were taken, not every possible temperature.
Peirce calls this an empirical formula when it describes the observed
relation without an explanatory hypothesis supplying the mechanism.
\end{hbexample}

The example contains no supplied dataset or measured coefficients.
Accordingly no numerical fit, uncertainty interval or actual experiment
is reported here. Peirce distinguishes error in observations from
inadequacy of the formula fitted to them. Even perfect observations
would not prove that a cubic form is exact. Replacing it by an infinite
expression would create another representation problem, not eliminate
the empirical issue.

His discussion of measurement scales also qualifies simplicity. A change
of units can require numerical constants even where the underlying relation
is unchanged. The contrast between an empirical expression and a physical
law is his methodological contrast in this paper, not a theorem that
counting fitted coefficients detects a true law.

The kinetic theory of gases supplies a different explanatory task.
Peirce describes a gas through moving particles and their impacts, aiming
to account for the approximate pressure--volume relation and its departures.
Finite particle size and attraction are offered as reasons for deviations.
His discussion combines the observed relation, a proposed microscopic
mechanism and assumptions carried over from mechanics. He notices that
applying macroscopic mechanical principles to unobserved molecules is
itself part of the hypothetical step
\parencite[sec.~IV]{Peirce1878DIHPrimary}.

The historical account of Bernoulli, heat and gas expansion is Peirce's
account, not a newly established chronology or current complete molecular
theory. It supplies no apparatus, population of trials, measurement table
or uncertainty analysis. Its inferential point survives that boundary:
deducing a consequence under a mechanism does not independently establish
that the mechanism applies. An account which also addresses departures
has more to explain than a formula reproducing the original regularity.

Peirce further observes that extending an induction far from its observed
range can acquire the character of hypothesis. Explanatory connections
then support an extension which resembles the original observations less
closely. The paper therefore does not insist that every scientific
reasoning episode belongs exclusively to one box.

\section{Why hypothesis is not merely induction awaiting replacement}
\label{sec:p16-peirce-1878-limits}

Section~V argues that hypotheses need not become inductions about the
same kinds of objects. Further documents can strengthen the conjecture
that Napoleon existed, but observing more documents does not turn an
observation of documents into direct observation of the earlier man.
The hypothesis can also have consequences of kinds not yet imagined.
Peirce's fanciful extraterrestrial photographs and reflected historical
light dramatize that possibility; they are not predicted observations
or established astronomy \parencite[sec.~V]{Peirce1878DIHPrimary}.

This gives a reason for distinguishing the inferential tasks without
claiming that every explanation must concern an unobservable cause.
Peirce also says induction is generally stronger. That comparative
assertion belongs to his account; the paper gives no common probability
scale ordering every possible induction above every possible hypothesis.
His own introduction includes a range from tentative surmise to familiar
historical belief.

The paper adds a psychological and physiological proposal: induction
forms a general habit, hypothesis unifies a tangle of predicates into
a conception and a corresponding feeling, and deduction directs
attention. Orchestral sound and its associated emotion illustrate the
proposed unity of feeling. The proposal is not a measured neural
mechanism. No recruited sample, task design, physiological recordings
or competing-model comparison is supplied.

Likewise, Peirce's division of classificatory, theoretical and hypothetical
sciences reflects his 1878 classification. It is not a claim that a
current discipline uses only one inference type. His footnote attributes
the division to his 1866 Lowell lectures and an 1867 publication. That
is his chronology statement; it does not provide the distinct complete
earlier presentations.

\section{The later vocabulary: argument, musement and retroduction}
\label{sec:p16-peirce-1908-context}

The 1908 article distinguishes an \emph{Argument}, a process tending
reasonably toward a definite belief, from an \emph{Argumentation},
which proceeds through explicitly formulated premises. Thus a reasoning
episode need not start with a finished list of propositions. That
distinction matters when the conjecture helps the inquirer articulate
what had initially been surprising.

Its opening defines the Real through identifying properties which belong
to their subject independently of whether a person or group attributes
them. The Actual concerns what is encountered in past, present or future.
His dream example separates the dream's imagined substance from the fact
that a dated dream occurred. The latter can be Real even when its content
does not describe an actual object. Peirce introduces
three Universes of Experience: Ideas, brute actuality, and the connecting
power exemplified by signs. These are philosophical commitments in the
article, not an empirical inventory proved exhaustive by naming three
categories. His flame-and-habit illustration concerns experience and
self-control; it is not advice to perform a hazardous experiment.

\emph{Musement} is his name for an open, playful exploration which can
notice wonders and suggest explanations without starting with a desired
verdict. He explicitly contrasts that attitude with setting out to
become convinced of religion. The latter threatens the candor which
inquiry requires. He permits diagrams and experiments in reflection,
while distinguishing such exploration from a systematic scientific
study \parencite[sec.~I]{Peirce1908NeglectedPrimary}.

His suggestion that human understanding can reach the causes of phenomena
is an encouragement to inquire, not a proof that every problem is soluble.
The text's conjectures about chemical elements, beauty and growth are
historical examples within its metaphysical argument, not present
scientific findings. His suggested times or proportions of waking life
for musement likewise have no supplied effectiveness study.

\section{Three stages, and the different work of testing}
\label{sec:p16-peirce-stages}

\begin{hbdefinition}{def:p16-peirce-retroduction}{Retroduction in the 1908 inquiry account}
The first stage begins with a phenomenon which disrupts an expectation.
It includes searching for pertinent circumstances, forming a conjecture,
recognizing its explanatory fit and provisionally appraising its
plausibility. An Explanation, in Peirce's defined sense, presents the
surprising fact as a consequence of the circumstances together with
the conjecture, if that conjecture is true.

Retroduction is the characteristic movement from a consequent toward
a proposed antecedent. It does not establish the antecedent's truth.
Peirce's first-stage acceptance ranges from a question worth examining
to a strong inclination to believe; the strength of that inclination
does not remove the need for testing.
\end{hbdefinition}

The first stage is therefore broader than one inversion of a finished
syllogism. The phenomenon may become articulable only in light of the
new conception. This explains why Peirce invokes Argument rather than
requiring Argumentation throughout. It also makes a reconstruction's
inputs consequential: a modern formalization which fixes all candidate
propositions in advance has not thereby captured how they were formed.

The \emph{second stage}, deduction, starts from the hypothesis rather
than collecting more appearances of the original surprising fact.
It first \emph{explicates} the hypothesis, making its content distinct,
then demonstrates conditional experiential consequences. Peirce relates
demonstration to diagrams or Icons, actually connecting Indices, and
interpreted Symbols. He distinguishes corollarial demonstration, which
uses considerations already involved in the explication, from a more
complicated theorematic process. He says the fuller definitions and
proofs cannot be given in this short article.

The purpose here is to obtain testable consequents. A correct conditional
deduction does not show that the hypothesis or its empirical auxiliary
premises are true. Clarifying a conjecture also does not manufacture
the experience needed to test its application.

The \emph{third stage}, induction, compares those consequents with
experience and judges whether the hypothesis is sufficiently correct,
needs modification or must be rejected. Peirce divides it into
classification of experiential objects, individual testing arguments
or \emph{Probations}, and a final appraisal of the tests separately,
together, and of the appraisals themselves. The method is not merely
counting supporting instances: it includes scrutinizing how the
assessment was obtained \parencite[sec.~III]{Peirce1908NeglectedPrimary}.

Crude induction, in this account, infers a universal negative from
the lack of an established counterinstance. Peirce emphasizes its
weakness through the former rejection of stones falling from the sky.
Gradual induction revises an estimate with further instances; he
distinguishes qualitative from quantitative forms, and measurements,
statistics and countings within the latter. This is a broader role
for induction than the earlier case--result-to-rule description alone.

He asserts a long-run correction property for gradual induction.
The article does not give a data-generating law, an error functional,
sampling conditions or the omitted proof needed to turn that assertion
into a mathematical convergence theorem. A contemporary consistency
theorem must state its own conditions. Peirce's claim cannot be credited
as their unrestricted historical proof.

\section{The argument for trying a conjecture, and its limits}
\label{sec:p16-peirce-validity}

Section~IV explicitly announces opinions whose complete positive
proofs are not supplied there. Peirce connects deduction with our ability
to explicate meaning and induction with a method capable, under continued
inquiry, of diminishing error. He assigns the introduction of new
conceptions to retroduction. Deduction clarifies and draws consequences;
induction appraises. Those claims describe the intended division of
labor, not three independently proved universal adequacy results
\parencite[secs.~III--IV]{Peirce1908NeglectedPrimary}.

His treatment of instinct contains an important internal objection.
The assertion that one cannot help accepting a conjecture is not, by
itself, the logical justification a rational agent owes. He expressly
rejects that excuse. He then proposes an analogy with specialized animal
capacities, combined with the historical success of scientific conjecture:
human minds may have a natural aptitude for grasping features of nature.

This is an abductive justification of the practice, not an experience-free
deduction of the truth of its outputs. The article appeals to computations
and histories which it does not print. Their announced existence supplies
neither a reproducible calculation nor a success-rate estimate. Its
animal analogy is not a controlled cognitive experiment. The additional
thought that abandoning every hypothesis would halt further knowledge
gives a practical reason to continue trying; it does not show that a
particular conjecture is right.

Peirce also corrects his own interpretation of simplicity. He says
he formerly took Galileo's preference to concern what is logically
simpler or adds least to observation. Against that reading he observes
that one might then refuse hypotheses altogether, merely list the
observations, and miss the unexpected complexity disclosed by scientific
advance. He instead emphasizes the hypothesis which comes more naturally
to a prepared inquirer. He does not say logical simplicity has no value;
he makes it secondary in this argument.

The proposed natural affinity remains a substantive and fallible
philosophical premise. It is not equivalent to minimum description
length, fewest assumed atoms, evolutionary fitness, a numerical prior,
or the strongest immediate feeling. A method's use, popularity or
practical necessity alone establishes none of those truth connections.

\section{The theological application exposes an unfinished warrant}
\label{sec:p16-peirce-theological-boundary}

The article's application must not be mistaken for a completed proof
that the first stage secures every desired belief. Its neglected argument
begins with musement about the three Universes, connectedness and growth,
leading toward the hypothesis of a necessary creator. Peirce distinguishes
that living course of thought from the theologian's formulated defense
and the further logical defense of the whole process.

Section~II admits a special difficulty: the proposed infinite object
exceeds a definite finite conception, yet a hypothesis must have some
conceived content. Peirce treats the conception as vague and progressively
developed. Purpose, which ordinarily involves growth, cannot simply be
attributed without qualification to the stipulated necessary being.
These are his interpretive qualifications, not a mathematical existence
proof obtained from a definition.

The relation between an attractive conception, practical commitment and
truth remains consequential. Peirce describes belief through preparedness
to shape conduct, and in Section~V assigns a role to self-controlled
growth under criticism. He still says the hypothesis needs Probation.
An inclination to organize one's life around it does not itself provide
an independently established experimental result
\parencite[secs.~II, V]{Peirce1908NeglectedPrimary}.

The paper's three imagined types of hearers are not a recruited population.
Near its close Peirce admits that little established fact supports the
claim that musement will convince every normal person. His reliance on
his own intellectual disposition, the prevalence of optimism and his
hostile classifications of pessimists does not supply a valid population
study or clinical diagnosis. Those classifications are not adopted here
as medical judgments or as grounds for dismissing disagreement.

He further acknowledges that the argument in his earlier popular essay
about the settlement of belief could be accused of begging the question.
He distinguishes actual satisfaction from the satisfaction which would
survive inquiry carried to its ultimate issue. That correction prevents
substituting present comfort for his proposed endpoint; it does not
prove that every inquiry reaches it. His polemic with James, Schiller
and other pragmatists is his stated contrast, not a complete independent
presentation of their positions.

\section{Reference use without silently modernizing the account}
\label{sec:p16-peirce-reference-use}

The 1878 and 1908 texts share a separation between an explanatory proposal
and deductive necessity, but their organizing tasks differ. The first
classifies synthetic conclusions through rule, case and result, allows
mixed scientific reasoning, and includes prediction safeguards. The
second places conjectural acceptance within a larger sequence which
must explicate consequences and appraise experience. It explicitly
recognizes that generating the conception and stating its premises are
different achievements. These observed differences do not amount to
a complete chronology of Peirce's revisions between the two publications.

The finite input of \cref{def:p16-abduction-input} already supplies
observations, background formulas and permitted atomic hypotheses.
Its exhaustive construction answers a precisely defined explanation
search. Peirce's first stage can instead concern what to consider as
a hypothesis at all. Moreover, the countervaluation in
\cref{thm:p16-abduction-truth-boundary} separates explanatory sufficiency
from the truth of an explanation. That is a justified neighboring
comparison, not an attribution of its model or proof to Peirce.

\begin{hbexample}{def:p16-peirce-inquiry-record}{Carrying the torn-paper conjecture into inquiry}
In a reconstruction of the source example, first record the writing,
the matching edges and the alleged access conditions. The authorship
conjecture is a proposed explanation, not an observation of authorship.
An alternative separates ownership of the sheet from who wrote on it.

Next explicate what would distinguish the conjectures: provenance of
the sheet, the opportunity for prior writing, and further independently
specified matches or records. Derive the proposed test consequences
under explicit auxiliary assumptions. Do not label an expected result
a deduction unless those assumptions actually connect it to the hypothesis.

Finally obtain and appraise the evidence, recording failures and the
selection of tests as well as successes. A change of explanation,
an auxiliary-premise revision and a measurement failure are different
possible responses. No tests were conducted in this illustrative
reconstruction, and no conclusion about an actual author follows.
\end{hbexample}

\begin{hbmethod}{def:p16-peirce-consultation}{Separate proposal, consequence and appraisal}
Identify the surprising observation and the background expectation.
State the proposed explanation and genuine alternatives. Distinguish
forming that proposal from granting it provisional attention, and both
from believing it established.

For testing, make the conjecture and auxiliary assumptions sufficiently
distinct to derive consequences. Specify how observations and test
families are selected; retain failures and possible error sources.
Appraise both the results and the procedure which produced the appraisal.

State what remains conditional or unknown. Explanatory fit is not
deductive truth; natural attraction is not measured reliability;
further testing is not an assurance of eventual success without further
conditions. A source's announced but unprinted proof remains unprinted,
even when its methodological suggestion is useful.
\end{hbmethod}

\index{Peirce, Charles Sanders!hypothesis and inquiry}
\index{abduction!Peirce's historical presentations}
\index{retroduction!stages of inquiry}
\index{hypothesis!proposal and testing}
\index{induction!Peirce's classification and appraisal}
\index{musement!Peirce}
